\documentclass[11pt]{article}

\usepackage[margin=1in]{geometry}
\usepackage{amsmath,amssymb,amsthm,mathtools}
\usepackage{booktabs}
\usepackage{array}
\usepackage{microtype}
\usepackage[hidelinks]{hyperref}

\newtheorem{theorem}{Theorem}
\newtheorem{proposition}[theorem]{Proposition}
\newtheorem{lemma}[theorem]{Lemma}
\newtheorem{remark}[theorem]{Remark}
\theoremstyle{definition}
\newtheorem{definition}[theorem]{Definition}
\theoremstyle{plain}
\newcommand{\R}{\mathbb{R}}
\newcommand{\C}{\mathbb{C}}
\newcommand{\Q}{\mathbb{Q}}
\newcommand{\Span}{\operatorname{span}}
\newcommand{\card}{\operatorname{card}}
\newcommand{\crs}{\operatorname{cr}}

\title{An Exact Real Quantum Latin Square of Order Six and Cardinality 29\\
\large Unsolved Labs Research Release R011}
\author{Unsolved Labs}
\date{Revised 18 August 2026}

\begin{document}
\maketitle

\begin{abstract}
We give an explicit real quantum Latin square of order six with cardinality \(29\).
The construction is defined over a quadratic number field
\(\Q(\tau)\), where
\[
123201\tau^2-202800\tau-93848=0
\]
and \(\tau\) is the positive root.
Thirty off-diagonal directions are constructed in \(\R^5\); after normalization they form orthonormal bases in every punctured row and column, and a common orthogonal diagonal vector extends the array to a quantum Latin square in \(\R^6\subset\C^6\).
Exact projective comparison shows that the thirty off-diagonal cells determine twenty-eight rays, so the common diagonal ray gives total cardinality \(29\).
The repository accompanying this manuscript contains a standard-library exact reconstruction, a frozen exact coordinate certificate checked by a smaller data-driven verifier, and a separately written SymPy reconstruction.
A distinct public cardinality-\(29\) construction by Das and Kumar appeared on arXiv before the R011 repository was initialized; accordingly, no priority claim is made here.
\end{abstract}

\section{Introduction}

A quantum Latin square (QLS) is a vector-valued analogue of an ordinary Latin square: its entries are unit vectors and every row and column is an orthonormal basis.
The notion was introduced by Musto and Vicary \cite{MustoVicary2016} and is connected to unitary error bases and quantum combinatorial design.

For order six, recent work has substantially clarified the attainable cardinalities, where cardinality means the number of entry rays after global phase is identified.
Xu's August 2026 revision \cite{Xu2026} left cardinality \(29\) as the remaining unresolved value in the order-six spectrum described there.
Das and Kumar subsequently posted an explicit real cardinality-\(29\) construction on 12 August 2026 \cite{DasKumar2026}.
The R011 repository was initialized on 14 August 2026, so the present artifact does not claim priority for resolving that gap.

The purpose of this release is instead to present and verify another explicit algebraic construction.
Its distinguishing public feature is a quadratic parameter \(\tau\) and a compact sequence of rational rotations, orthogonal projections, and generalized cross products.
The proof is finite and exact.
The repository freezes the reduced coordinates \(a+b\tau\) of all thirty off-diagonal directions and checks the full QLS and cardinality conditions without floating-point arithmetic.

\paragraph{Main result.}
\begin{theorem}\label{thm:main}
There exists a real quantum Latin square of order \(6\) with cardinality \(29\).
\end{theorem}

\paragraph{Verification architecture.}
The human proof below is accompanied by three exact computational views of the same finite theorem:
(i) a construction-level verifier over \(\Q(\tau)\) implemented using only Python's standard library;
(ii) a frozen coordinate certificate checked by a smaller program that does not import the construction;
and (iii) a separately written SymPy reconstruction.
The search or discovery process is not part of the proof trust boundary.

\section{Definitions and reduction to a punctured array}

\begin{definition}
A quantum Latin square of order \(n\) is an array
\(\Phi=(\phi_{ij})_{0\leq i,j<n}\) of unit vectors in \(\C^n\) such that each row and each column is an orthonormal basis of \(\C^n\).
Its cardinality is the number of rays represented by its entries, where two nonzero vectors are identified when one is a nonzero complex scalar multiple of the other.
\end{definition}

We use a common diagonal vector and work first in a five-dimensional orthogonal complement.

\begin{lemma}[Diagonal extension]\label{lem:diagonal}
Let \(V=(v_{ij})_{i\neq j}\) be an order-\(6\) punctured array of nonzero vectors in \(\R^5\) such that, after normalizing the vectors, every punctured row and every punctured column is an orthonormal basis of \(\R^5\).
Let
\[
f=(0,0,0,0,0,1)^{\mathsf T}\in\R^6
\]
and, for \(i\neq j\), set
\[
\phi_{ij}=\frac{(v_{ij},0)}{\|v_{ij}\|},\qquad \phi_{ii}=f.
\]
Then \(\Phi=(\phi_{ij})\) is a real QLS(6).
If the off-diagonal entries determine \(r\) rays in \(\R^5\), then \(\Phi\) has cardinality \(r+1\).
\end{lemma}

\begin{proof}
The common diagonal vector is orthogonal to the embedded copy of \(\R^5\).
Thus each normalized punctured row or column, together with \(f\), is an orthonormal basis of \(\R^6\).
The ray of \(f\) is disjoint from every ray in the embedded \(\R^5\), giving the cardinality formula.
\end{proof}

\section{Quadratic parameter and building blocks}

Let
\[
P(T)=123201T^2-202800T-93848,
\]
and let
\[
\tau=\frac{2600+2\sqrt{3590422}}{3159}
\]
be its positive root.
We work in \(\Q(\tau)\subset\R\).

Let \(e_0,\ldots,e_4\) be the standard orthonormal basis of \(\R^5\).
Define
\[
a=\frac{4e_3+3e_4}{5},\qquad
p=\frac{-3e_3+4e_4}{5},
\]
\[
d=\frac{5e_1-12e_2}{13},\qquad
t=\frac{12e_1+5e_2}{13},
\]
and
\[
q=\frac{8t+15e_4}{17},\qquad
r=\frac{15t-8e_4}{17}.
\]

Split \(q=X+H\) by
\[
X=\frac{96}{221}e_1+\frac{12}{17}p,\qquad
H=\frac{40}{221}e_2+\frac{9}{17}a,
\]
and introduce
\[
Y=\frac{12}{17}e_1-\frac{96}{221}p,\qquad
U=\frac{9}{17}e_2-\frac{40}{221}a.
\]
A direct calculation gives
\[
\sigma=\langle X,X\rangle=\frac{33552}{48841},\qquad
\rho=\langle H,H\rangle=\frac{15289}{48841},\qquad
\sigma+\rho=1.
\]
Set
\[
S=\rho X-\sigma H,
\]
\[
W=32e_0-51q,\qquad Z=51e_0+32q,
\]
and
\[
B=51\rho e_0+32H,\qquad C=32e_0-51H.
\]

Now put
\[
T_1=4e_0-3t,\qquad T_2=d+\tau e_3.
\]
The vectors \(e_4,T_1,T_2\) are mutually orthogonal and have squared norms \(1,25,1+\tau^2\).
Define
\[
G=X-\langle X,e_4\rangle e_4
-\frac{\langle X,T_1\rangle}{25}T_1
-\frac{\langle X,T_2\rangle}{1+\tau^2}T_2.
\]
Thus \(G\) is the orthogonal projection of \(X\) to
\(\Span(e_4,T_1,T_2)^\perp\).

Define
\[
K=(600\tau+1521)e_0+1872e_1+(2080\tau+780)e_2+1920e_3.
\]
Exact reduction modulo \(P(\tau)=0\) gives
\[
K\perp e_4,\qquad K\perp T_1,\qquad K\perp T_2,\qquad K\perp X.
\]
Consequently \(K\perp G\).

Finally let
\[
g_U=\langle G,U\rangle,\qquad g_Y=\langle G,Y\rangle,
\]
\[
L=g_UY-g_YU,\qquad
M=g_Y\rho Y+g_U\sigma U.
\]

For four vectors \(v_1,\ldots,v_4\in\R^5\), let
\(\crs(v_1,v_2,v_3,v_4)\) denote the generalized cross product whose
\(j\)-th coordinate is \((-1)^j\) times the determinant obtained by deleting column \(j\) from the \(4\times5\) matrix with rows \(v_i^{\mathsf T}\).
Set
\[
N_3=\crs(L,B,W,G),\qquad N_4=\crs(M,C,X,K).
\]

\section{The punctured array}

Define the off-diagonal directions by
\[
V=
\begin{pmatrix}
- & e_0 & e_1 & e_2 & e_3 & e_4\\
3e_0+4t & - & -p & a & d & T_1\\
\tau d-e_3 & q & - & e_0 & r & T_2\\
N_3 & L & B & - & W & G\\
N_4 & M & C & X & - & K\\
W & S & U & Y & Z & -
\end{pmatrix}.
\]
All thirty displayed off-diagonal directions are nonzero in \(\Q(\tau)^5\).
The repository freezes their reduced coordinates in the basis \((1,\tau)\) as a JSON certificate.

\section{Orthogonality}

The first three rows follow from the rational rotations above.
The decomposition \(q=X+H\), together with the orthogonal companion pairs \((X,Y)\) and \((H,U)\), supplies the relations needed for row \(5\) and columns \(1\) through \(4\).

Column \(5\) is
\[
e_4,\ T_1,\ T_2,\ G,\ K.
\]
By construction \(G\perp e_4,T_1,T_2\), and \(K\perp e_4,T_1,T_2,X\).
Since \(G\) is the projection of \(X\) into their common orthogonal complement, \(K\perp G\).

The remaining scalar closures are exactly where the quadratic parameter enters:
\[
\langle B,G\rangle
=
\frac{4608P(\tau)}
     {206353225(1+\tau^2)}
=0,
\]
and
\[
\langle M,K\rangle
=
-\frac{154607616(200\tau+507)P(\tau)}
       {171334463657825(1+\tau^2)}
=0.
\]
The definitions of \(L,M,N_3,N_4\) force the remaining row-\(3\) and row-\(4\) products to vanish.

For completeness, one may recover column \(0\) by a frame identity.
After normalization, the five off-diagonal vectors in each row form an orthonormal basis of \(\R^5\), so the sum of all row projectors is \(6I_5\).
Columns \(1,\ldots,5\) contribute \(5I_5\), leaving \(I_5\) as the frame operator of column \(0\).
Five vectors in dimension five with frame operator \(I_5\) form an orthonormal basis.

The exact verifiers do not rely on omitting this column: they evaluate all \(120\) pairwise row and column orthogonality equations directly.

\begin{proposition}\label{prop:punctured}
After normalization, every punctured row and every punctured column of \(V\) is an orthonormal basis of \(\R^5\).
\end{proposition}

\begin{proof}
The preceding construction proves the structural relations.
The finite remainder is the collection of exact inner products described above; the repository's exact verifiers reduce every one of them to zero in \(\Q(\tau)\).
\end{proof}

\section{Ray count}

Two nonzero vectors \(x,y\in\Q(\tau)^5\) determine the same ray if and only if all \(2\times2\) minors
\[
x_i y_j-x_j y_i
\]
vanish.
The release evaluates this criterion for all
\[
\binom{30}{2}=435
\]
pairs of off-diagonal cells.

Exactly two off-diagonal ray classes repeat:
\[
V_{01}\sim V_{23},\qquad V_{34}\sim V_{50}.
\]
Indeed the first equality is literally \(e_0=e_0\), while the second is literally \(W=W\).
No other pair is proportional.

A phase-class labeling of the off-diagonal entries is
\[
\begin{pmatrix}
D&0&1&2&3&4\\
5&D&6&7&8&9\\
10&11&D&0&12&13\\
14&15&16&D&17&18\\
19&20&21&22&D&23\\
17&24&25&26&27&D
\end{pmatrix},
\]
where \(D\) denotes the common diagonal ray.
Thus the thirty off-diagonal cells determine twenty-eight rays, and the diagonal contributes one new ray.

\begin{proof}[Proof of Theorem~\ref{thm:main}]
By Proposition~\ref{prop:punctured} and Lemma~\ref{lem:diagonal}, the normalized array is a real QLS(6).
The exact projective classification gives twenty-eight off-diagonal rays and the common diagonal ray is distinct from all of them.
Therefore
\[
\card(\Phi)=28+1=29.
\]
\end{proof}

\section{Exact verification and trust boundary}

The accompanying repository provides three complementary exact checks.

\paragraph{Construction verifier.}
\texttt{verify\_qls6\_card29.py} reconstructs the vectors from the displayed formulas using only Python's standard library and exact \texttt{Fraction} arithmetic in the basis \((1,\tau)\).

\paragraph{Frozen certificate checker.}
\texttt{certificate/qls6\_card29\_exact.json} freezes all thirty exact directions as rational coefficient pairs \(a+b\tau\).
The smaller \texttt{verify\_certificate.py} does not import the construction code; it checks nonzeroness, all \(120\) orthogonality identities, all \(435\) projective comparisons, the two repeated classes, and the final cardinality directly from the certificate.
A separate deterministic bridge regenerates the certificate from the construction and demands byte-for-byte equality.

\paragraph{Independent symbolic reconstruction.}
\texttt{verify\_qls6\_card29\_sympy.py} separately implements the formulas in SymPy and reduces rational functions modulo \(P(T)\).

No decimal approximation is used as evidence for the theorem.
The decimal vector export is retained only for inspection.

\paragraph{Formalization boundary.}
This release does not claim Lean verification.
The finite theorem has instead been reduced to an explicit exact coordinate certificate checked by a deliberately small standard-library program.
A proof-assistant claim would require a faithful encoding of the quadratic field, all thirty directions, the orthogonality/projective checks, and the final QLS reduction; such a formalization is not part of R011.

\section{Related work and scope}

Musto and Vicary introduced quantum Latin squares \cite{MustoVicary2016}.
Recent papers have developed the cardinality problem, including the order-six constructions of Xu \cite{Xu2026} and broader spectrum results such as \cite{ZhangLvCao2026}.

Das and Kumar \cite{DasKumar2026} publicly posted another real cardinality-\(29\) QLS before the R011 repository was initialized.
Their public construction uses a different displayed parameterization based on rational unit-circle parameters and orthogonal changes of basis.
The present manuscript therefore makes no claim of first discovery or first closure of the order-six spectrum.

The result established here is intentionally narrow:
one explicit real QLS(6) has cardinality \(29\).
No uniqueness, classification, or coordinate-optimality result is asserted.

\section*{Research provenance}

This manuscript and repository form an Unsolved Labs public research artifact produced in a frontier-AI research workflow.
No conventional human-authorship claim is implied by the organizational author line.
Correctness is supported by the explicit mathematical argument and exact machine-verification artifacts, not by the provenance of the discovery process.

\begin{thebibliography}{99}

\bibitem{MustoVicary2016}
B.~Musto and J.~Vicary,
\emph{Quantum Latin squares and unitary error bases},
Quantum Information and Computation 16 (2016), 1318--1332.
doi:10.26421/QIC16.15-16-4; arXiv:1504.02715.

\bibitem{Xu2026}
Z.~Xu,
\emph{New Cardinalities for Quantum Latin Squares of Order Six},
arXiv:2607.11800v3 (revised 4 August 2026).

\bibitem{DasKumar2026}
A.~P. Das and D.~Kumar,
\emph{A Quantum Latin Square of Order Six with Cardinality 29},
arXiv:2608.12607 (submitted 12 August 2026).

\bibitem{ZhangLvCao2026}
Y.~Zhang, M.~Lv, and H.~Cao,
\emph{On the possible cardinalities of quantum Latin squares},
arXiv:2607.19969 (2026).

\end{thebibliography}

\end{document}
